% =============================================================
% บทที่ 5 - สรุปผลและอภิปรายผล
% โครงงาน: การออกแบบและพัฒนาเกมเดินเรือด้วยระบบสั่งการด้วยเสียงเพื่อส่งเสริมทักษะการสื่อสาร
% =============================================================

\section{ผลการดำเนินงานที่แล้วเสร็จ}

รายงานความก้าวหน้าฉบับนี้สรุปผลการดำเนินงานของโครงงานการออกแบบและพัฒนาเกมเดินเรือด้วยระบบสั่งการด้วยเสียงเพื่อส่งเสริมทักษะการสื่อสาร ณ วันที่ 9 ตุลาคม 2569 ซึ่งครอบคลุมการศึกษาทฤษฎีและงานวิจัยที่เกี่ยวข้อง การออกแบบ การพัฒนา และการทดสอบระบบ โดยมีรายละเอียดดังนี้

ในด้านการออกแบบ ผู้จัดทำได้กำหนดเนื้อเรื่องและสตอรี่บอร์ด กระบวนการเล่นหลัก (Core Game Loop) สถาปัตยกรรมของระบบ ส่วนติดต่อผู้ใช้ (UI/HUD) และบทบาทของตัวละครทั้งสามตำแหน่ง ได้แก่ กัปตัน (ผู้เล่น) ต้นกลเรือ และพลถือท้ายและยามระวังหน้า รวมถึงจัดทำคลังคำสั่งการเดินเรือมาตรฐาน (Standard Marine Communication Phrases: SMCP) สำหรับใช้เป็นเกณฑ์ตรวจสอบความถูกต้องของคำสั่งเสียง ดังแสดงในภาคผนวก ง

ในด้านการพัฒนา ผู้จัดทำได้พัฒนาระบบหลักของเกมบน Unity ด้วยภาษา C\# ครบทุกระบบ ได้แก่ (1) ระบบวิทยุรับคำสั่งเสียง ซึ่งแปลงเสียงเป็นข้อความด้วยโมเดล ggml-base.en ผ่าน whisper.cpp ภายในเครื่องหลังผู้เล่นปล่อยปุ่มวิทยุ (2) ตัวแยกคำสั่ง (Command Parser) ซึ่งตีความข้อความภายในเครื่องก่อน และส่งข้อความไปยังแบบจำลองภาษา GPT-4o mini ผ่าน OpenAI API เฉพาะเมื่อไม่พบคำสั่งที่ชัดเจนหรือเมื่อผู้เล่นสนทนากับลูกเรือปัญญาประดิษฐ์ (AI Crew) (3) ระบบฟิสิกส์ของเรือซึ่งจำลองแรงเฉื่อยและรัศมีวงเลี้ยวของเรือขนาดใหญ่ (4) ระบบสุ่มสร้างสิ่งกีดขวาง (5) ส่วนแสดงผลบนหน้าจอ (HUD) (6) บทฝึกสอน (7) โหมดเนื้อเรื่องแบบ 3 บท บทละ 3 ภารกิจ (8) เครื่องมือช่วยสอนภาษา เช่น คำแนะนำประโยคและการทวนคำสั่งด้วยรูปประโยคมาตรฐาน (9) อู่ต่อเรือและระบบอัปเกรด และ (10) ระบบบันทึกข้อมูลสำหรับการวิจัย (turns.csv, voyages.csv และ emergencies.csv)

ในด้านการทดสอบ ชุดทดสอบหน่วย (Unit Test) จำนวน 363 กรณีผ่านทั้งหมดดังตารางที่ \ref{tab:unit_result} ส่วนกรณีทดสอบระดับระบบมีผลตามตารางที่ \ref{tab:test_result}--\ref{tab:test_result2} และการทดลองเล่นของผู้จัดทำเมื่อวันที่ 8 ตุลาคม 2569 ให้ค่าเบื้องต้นของเวลาแปลงเสียงเป็นข้อความอยู่ระหว่าง 2.47--2.87 วินาที (n = 7) และเรือเริ่มทำตามคำสั่งภายใน 2.47 และ 2.56 วินาที (n = 2) ดังตารางที่ \ref{tab:nfr_result} ทั้งนี้เป็นการวัดเบื้องต้นจากผู้จัดทำเพียงคนเดียว ยังไม่ใช่ผลการประเมินกับกลุ่มตัวอย่าง

งานที่ยังคงเหลือ ได้แก่ การแก้ไขกรณีทดสอบที่ผ่านบางส่วนในตารางที่ \ref{tab:test_result}--\ref{tab:test_result2} การทดสอบกรณีที่ไม่มีคีย์ API และกรณีที่การเชื่อมต่อ OpenAI API หมดเวลา (Timeout) และการประเมินผลกับกลุ่มตัวอย่างจำนวน 10--20 คน ซึ่งยังไม่ได้เก็บข้อมูล โดยจะวัดความวิตกกังวลในการพูดภาษาอังกฤษก่อนและหลังเล่นด้วยแบบวัดที่ดัดแปลงจาก FLCAS~\cite{horwitz1986} ร่วมกับแบบสอบถามประสบการณ์การเล่นเกม และวัดความแม่นยำของระบบจากคำสั่งเป้าหมาย 10 คำสั่งต่อผู้เข้าร่วม ผลการประเมินดังกล่าวจะใช้สรุปว่าเกมช่วยลดความวิตกกังวลในการพูดภาษาอังกฤษได้มากน้อยเพียงใด

\section{ปัญหาและข้อจำกัด}

ในการดำเนินโครงงานและพัฒนาตัวเกมต้นแบบ ผู้จัดทำพบปัญหาและข้อจำกัดที่ส่งผลต่อประสิทธิภาพของระบบและการประเมินผล ดังนี้

\begin{enumerate}[label=\arabic*)]
    \item \textbf{ข้อจำกัดด้านฮาร์ดแวร์และสภาพแวดล้อม:} ความแม่นยำในการรับคำสั่งขึ้นอยู่กับคุณภาพของไมโครโฟนและเสียงรบกวนภายนอก (Background Noise) ขณะเล่น แม้เกมจะมีการตัดเสียงที่เบาเกินไปด้วยค่ารากกำลังสองเฉลี่ย (RMS) และกรองข้อความที่เกิดจากเสียงรบกวนก่อนส่งต่อแล้ว แต่ในห้องที่มีเสียงดังระบบยังอาจถอดความผิดพลาดได้
    \item \textbf{ข้อจำกัดด้านระบบการรู้จำเสียง:} โมเดล ggml-base.en ยังถอดความสำเนียง (Accent) ของผู้เล่นที่ไม่ใช่เจ้าของภาษาผิดพลาดได้ ตัวแยกคำสั่งจึงมีตารางคำที่ระบบมักได้ยินผิดและการจับคู่คำแบบใกล้เคียง (Fuzzy Matching) เพื่อชดเชย แต่ตารางดังกล่าวสร้างจากการทดลองเล่นของผู้จัดทำ ยังไม่ครอบคลุมสำเนียงของกลุ่มตัวอย่าง
    \item \textbf{ข้อจำกัดด้านขอบเขตการสื่อสาร:} การสนทนาแบบอิสระกับลูกเรือผ่าน GPT-4o mini ยังไม่ถูกประเมินผลเชิงการสอน คำสั่งที่ใช้ควบคุมเรือจำกัดเฉพาะวลีสั่งการมาตรฐาน (SMCP) ที่กำหนด และการสนทนาต้องเชื่อมต่ออินเทอร์เน็ต
    \item \textbf{ข้อจำกัดด้านการพึ่งพาบริการภายนอก:} การตีความข้อความที่ไม่ชัดเจนและการสนทนากับลูกเรือขึ้นอยู่กับ OpenAI API ซึ่งต้องใช้เครือข่าย มีค่าใช้จ่ายตามปริมาณการใช้งาน และระบบต้องส่งข้อความที่ถอดจากเสียงของผู้เล่นออกจากเครื่องคอมพิวเตอร์ไปประมวลผลภายนอก (ไม่มีการส่งหรือบันทึกเสียง)
    \item \textbf{ข้อจำกัดด้านแพลตฟอร์ม:} เสียงพูดของลูกเรือใช้ระบบสังเคราะห์เสียงของ Windows จึงรองรับเฉพาะระบบปฏิบัติการ Windows
    \item \textbf{ข้อจำกัดด้านกลุ่มตัวอย่างและการประเมินผล:} การประเมินผลใช้แบบแผนกลุ่มเดียววัดก่อนและหลังโดยไม่มีกลุ่มควบคุม กับกลุ่มตัวอย่าง 10--20 คนที่เลือกแบบอาสาสมัครภายในมหาวิทยาลัย ซึ่งเป็นขนาดของการศึกษานำร่อง จึงตรวจพบได้เฉพาะผลขนาดใหญ่ ไม่สามารถสรุปความเป็นเหตุเป็นผลได้อย่างเด็ดขาด และอาจไม่สะท้อนประชากรในวงกว้าง นอกจากนี้การทดลองมีเพียงครั้งเดียวประมาณ 30 นาที ผลที่ได้จึงสะท้อนความวิตกกังวลในการพูดภาษาอังกฤษระยะสั้นขณะทดลอง (State Anxiety) ไม่ใช่การเปลี่ยนแปลงในระยะยาว และเกมเป็นเพียงเครื่องมือเสริม ไม่สามารถใช้ทดแทนการทดสอบมาตรฐานทางภาษาได้
\end{enumerate}

\section{แนวทางการแก้ไขและข้อเสนอแนะ}

เพื่อให้การพัฒนาเกมมีความสมบูรณ์และสามารถนำไปต่อยอดได้ในอนาคต ผู้จัดทำมีแนวทางการแก้ไขและข้อเสนอแนะ ดังนี้

\begin{enumerate}[label=\arabic*)]
    \item \textbf{การจัดการเสียงรบกวน:} ควรแนะนำให้ผู้เล่นใช้ชุดหูฟังที่มีไมโครโฟนในตัวและเล่นในห้องที่เงียบ และในระยะถัดไปควรเพิ่มการลดเสียงรบกวนก่อนส่งเสียงเข้าสู่ whisper.cpp เช่น RNNoise
    \item \textbf{การปรับปรุงความแม่นยำในการรับคำสั่ง:} ควรขยายตารางคำที่ระบบได้ยินผิดจากข้อมูล turns.csv ของกลุ่มตัวอย่าง เพื่อรองรับสำเนียงที่หลากหลายยิ่งขึ้น และทดสอบการประมวลผลเสียงด้วยการ์ดจอบนเครื่องอื่นเพิ่มเติม โดยเฉพาะเครื่องที่ไม่มีการ์ดจอแยก (บนเครื่องทดสอบ การ์ดจอลดเวลาแปลงเสียงจากประมาณ 1.6 วินาทีเหลือประมาณ 0.06 วินาที)
    \item \textbf{การขยายขอบเขตในอนาคต:} ควรขยายคลังคำสั่ง (Command Vocabulary) ให้รองรับรูปแบบการสื่อสารทางทะเลที่ซับซ้อนยิ่งขึ้น ประเมินผลเชิงการสอนของการสนทนากับลูกเรือ และใช้แบบแผนที่มีกลุ่มควบคุมพร้อมการวัดซ้ำหลังผ่านไประยะหนึ่ง (Delayed Post-test) กับกลุ่มตัวอย่างที่ใหญ่ขึ้น เพื่อประเมินผลต่อความวิตกกังวลในการพูดภาษาอังกฤษในระยะยาว
\end{enumerate}